% XeLaTeX; copy into a trusted batch. Defaults: body 12pt / baseline 18pt.
\documentclass[UTF8,fontset=none,12pt,a4paper]{ctexart}
\usepackage[left=20mm,right=20mm,top=24mm,bottom=20mm,headheight=15pt,headsep=7mm]{geometry}
\usepackage{fontspec}
\IfFontExistsTF{SimSun}{
  \IfFontExistsTF{SimHei}{\setCJKmainfont{SimSun}[BoldFont=SimHei]}{\setCJKmainfont{SimSun}[AutoFakeBold=1.5]}
}{
  \IfFontExistsTF{Noto Serif CJK SC}{\setCJKmainfont{Noto Serif CJK SC}}{\setCJKmainfont{FandolSong-Regular}[BoldFont=FandolSong-Bold]}
}
\IfFontExistsTF{SimHei}{\setCJKsansfont{SimHei}[BoldFont=SimHei]}{
  \IfFontExistsTF{Noto Sans CJK SC}{\setCJKsansfont{Noto Sans CJK SC}}{\setCJKsansfont{FandolHei-Regular}[BoldFont=FandolHei-Bold]}
}
\IfFontExistsTF{Times New Roman}{\setmainfont{Times New Roman}\setsansfont{Times New Roman}}{\setmainfont{TeX Gyre Termes}\setsansfont{TeX Gyre Termes}}
\usepackage{amsmath,amssymb,array,graphicx,fancyhdr,needspace}
\usepackage[math-style=ISO,bold-style=ISO]{unicode-math}
\IfFontExistsTF{Cambria Math}{\setmathfont{Cambria Math}}{
  \IfFontExistsTF{TeX Gyre Termes Math}{\setmathfont{TeX Gyre Termes Math}}{\setmathfont{Latin Modern Math}}
}
\usepackage{enumitem,tikz,booktabs,tabularx}
\setlist[enumerate]{nosep,leftmargin=2em,topsep=3pt}
\newcommand{\bodyfont}{\rmfamily\fontsize{12}{18}\selectfont}
\AtBeginDocument{\bodyfont\setlength{\abovedisplayskip}{6pt plus 1pt minus 1pt}\setlength{\belowdisplayskip}{6pt plus 1pt minus 1pt}\setlength{\abovedisplayshortskip}{4pt}\setlength{\belowdisplayshortskip}{4pt}}
\setlength{\parindent}{0pt}
\setlength{\parskip}{4pt}
\setlength{\emergencystretch}{1em}
\pagestyle{fancy}
\fancyhf{}
\fancyhead[L]{{\rmfamily\fontsize{10.5}{13}\selectfont 理学优课高中部}}
\fancyhead[R]{{\rmfamily\fontsize{10.5}{13}\selectfont 针对补充训练}}
\fancyfoot[C]{{\rmfamily\fontsize{10.5}{13}\selectfont\thepage}}
\renewcommand{\headrulewidth}{0.4pt}
\newcommand{\doctitle}[2]{\begin{center}{\sffamily\bfseries\fontsize{18}{23}\selectfont #1}\par\vspace{1mm}{\bodyfont 姓名：#2}\end{center}\vspace{1mm}}
\newcommand{\compacttitle}[1]{\begin{center}{\sffamily\bfseries\fontsize{18}{23}\selectfont #1}\end{center}}
\newcommand{\parttitle}[1]{\par\Needspace{6\baselineskip}\vspace{2mm}{\sffamily\bfseries\fontsize{14}{19}\selectfont #1}\par\vspace{1mm}}
\newcommand{\answerheading}[1]{\par\Needspace{6\baselineskip}{\sffamily\bfseries\fontsize{12}{18}\selectfont #1}\par}
\newcommand{\question}[2]{\par\Needspace{5\baselineskip}\noindent\textbf{#1.}\quad #2\par}
% Number each subquestion on its own line and align any continuation lines.
\newcommand{\subquestion}[2]{%
  \par\Needspace{2\baselineskip}\noindent%
  \hangindent=2em\hangafter=1%
  \makebox[2em][l]{(#1)}#2\par%
}
\newcommand{\workarea}[1]{\par\vspace*{#1}\par}
% Keep an ordinary question and its writing area together. Split very long tasks manually.
\newbox\handoutbox
\newcommand{\handoutblock}[1]{%
  \par\setbox\handoutbox=\vbox{\hsize=\linewidth #1}%
  \ifdim\dimexpr\ht\handoutbox+\dp\handoutbox\relax>\textheight
    \PackageError{math-handout}{Question block exceeds one page}{Split this task at a subquestion; keep body font at 12pt.}%
  \fi
  \Needspace{\dimexpr\ht\handoutbox+\dp\handoutbox+6pt\relax}%
  \box\handoutbox\par
}
\newcommand{\questionwithspace}[3]{\handoutblock{\noindent\textbf{#1.}\quad #2\par\vspace*{#3}}}
\newcommand{\originalquestion}[3]{\handoutblock{{\sffamily\bfseries #1}\par #2\par #3\par}}
\newcommand{\keypoints}[1]{%
  \par\addvspace{6pt}%
  \handoutblock{\noindent\textbf{要点}\quad #1\par}%
  \addvspace{6pt}%
}
\raggedbottom
% The solution can flow across pages. Reserve room for each new subpart.
\newcommand{\solutionblock}[2]{\answerheading{#1}#2\par}
\newcommand{\solutionpart}[1]{%
  \par\Needspace{4\baselineskip}\addvspace{5pt}%
  \noindent\textbf{(#1)}\par\nobreak\ignorespaces%
}
